

\chapter{Thick-thin decomposition} \label{thick-thin:chapter}
A peculiar aspect of hyperbolic geometry is the existence of complete hyperbolic manifolds that have finite volume but are not compact. These manifolds behave like the compact ones in many aspects, but are sometimes easier to construct. They arise naturally when we study hyperbolic surfaces, and are a fundamental constituent of Thurston's geometrisation of three-manifolds.

We prove here a structure theorem for all such manifolds. The theorem says that every finite-volume complete hyperbolic manifold can be decomposed into two domains: a \emph{thick part} which is compact and has injectivity radius bounded from below, and a \emph{thin part} that consists of \emph{cusps}. A cusp is a (truncated) hyperbolic manifold of type $N\times [0,+\infty)$, where every section $N\times t$ has a flat metric that shrinks exponentially with $t$. 

The core of this theorem is a lemma about discrete subgroups of Lie groups called the \emph{Margulis Lemma}. We will apply this lemma also to the other geometries and prove Bierbach's Theorem, that states that every compact flat manifold is covered by a torus.

Throughout the discussion we will also study some general aspects of finite-volume complete hyperbolic manifolds, concerning in particular closed geodesics, isometry groups, and finite covers.

\section{Tubes and cusps} \label{tubes:cusps:section}
We introduce here two very simple kinds of hyperbolic manifolds, so simple that their fundamental groups will be called \emph{elementary} in the sequel: the tubes and the cusps. Before introducing them we show that the injectivity radius of a hyperbolic manifold behaves nicely.

We also study the closed geodesics in a hyperbolic manifold $M$, and prove that there is precisely one in every free homotopy class of closed curves of hyperbolic type.

\subsection{Injectivity radius}
Like many other geometric properties, the injectivity radius of a complete hyperbolic manifold may be nicely observed by looking at its universal cover. If $S\subset \matH^n$ is a discrete set, we define $d(S)$ as the infimum of $ d(x_1,x_2)$ among all pairs $x_1, x_2$ of distinct points in $S$. 

Recall that every complete hyperbolic manifold is a quotient $M=\matH^n/_\Gamma$ by some group $\Gamma$ of isometries acting freely on $\matH^n$. 

\begin{prop} \label{d:prop}
Let $M=\matH^n/_\Gamma$ be a complete hyperbolic manifold and $\pi\colon \matH^n \to M$ the projection. For every $x\in M$ we have
$$\inj_xM = \frac 12 \cdot d(\pi^{-1}(x)).$$
\end{prop}
\begin{proof}
The number $\inj_x M$ is the supremum of all $r>0$ such that $B(x,r)$ is isometric to a ball of radius $r$ in $\matH^n$. The open set $B(x,r)$ is a ball of radius $r$ if and only if its counterimage via $\pi$ consists of disjoint balls of radius $r$, and this holds $\Leftrightarrow$ two distinct points in $\pi^{-1}(x)$ stay at distance at least $2r$.
\end{proof}

Recall that we denote by $d(\gamma)$ the minimum displacement of $\gamma \in \Iso(\matH^n)$.

\begin{cor} \label{inj:cor}
Let $M=\matH^n/_\Gamma$ be a complete hyperbolic manifold. Then
$$\inj M = \frac 12 \cdot \inf \big\{d(\gamma)\ \big|\ \gamma \in \Gamma, \gamma \neq \id\big\}.$$
\end{cor}
\begin{proof}
We have
$$\inj_xM = \frac 12 \cdot d(\pi^{-1}(x)) = \frac 12 \cdot \inf \big\{ d(\tilde x, \gamma(\tilde x)) \ \big|\ \gamma\in \Gamma, \gamma\neq \id, \tilde x \in \pi^{-1}(x)\big\}. $$
Therefore
$$\inj M = \inf_{x\in M} \inj_x M = 
\frac 12 \cdot \inf \big\{ d(\gamma)\ \big| \ \gamma \in \Gamma, \gamma \neq \id \big\}.$$
The proof is complete.
\end{proof}

\begin{cor}
If $M=\matH^n/_\Gamma$ is a compact hyperbolic manifold then every non-trivial element in $\Gamma$ is hyperbolic.
\end{cor}
\begin{proof}
Every nontrivial element in $\Gamma$ is either hyperbolic or parabolic.
If $M$ is compact then $\inj M >0$. If $\Gamma$ contains a parabolic $\gamma$ then $d(\gamma)=0$ and hence $\inj M =0$.
\end{proof}

\subsection{Tubes} \label{tubes:subsection}
Consider the infinite cyclic group $\Gamma = \langle \varphi \rangle$ generated by a hyperbolic transformation $\varphi$ on $\matH^n$ with axis $l$ and minimum displacement $d>0$. The iterates $\varphi^k$ are again hyperbolic transformations with axis $l$ and displacement $kd$. Therefore $\Gamma$ acts freely on $\matH^n$.
The quotient manifold $M=\matH^n/_\Gamma$ is called an \emph{infinite tube}. \index{tube}

\begin{ex}
Fix $q\in l$. Let $q_1,q_2$ be the two points in $l$ at distance $\frac d2$ from $q$ and $\pi_1,\pi_2$ the two hyperplanes orthogonal to $l$ in $q_1,q_2$.
The Dirichlet domain $D(q)$ of $\Gamma$ is the space comprised between $\pi_1$ and $\pi_2$.
\end{ex}

The infinite tube $M = \matH^n/_\Gamma$ is obtained from $D(q)$ by identifying $\pi_1$ and $\pi_2$ along $\varphi$. Its fundamental group is isomorphic to $\Gamma \Isom \matZ$. The axis $l$ projects in $M$ onto a closed geodesic $\gamma$ of length $d$. We have $\inj M = \frac d2 $ by Corollary \ref{inj:cor} and the points in $\gamma$ are precisely those with minimum injectivity radius. The closed geodesic $\gamma$ is sometimes called the \emph{core geodesic} of the tube.

\begin{prop}
Every infinite tube is diffeomorphic to $S^1\times \matR^{n-1}$ or $S^1 \timtil \matR^{n-1}$ according to whether $\varphi$ is orientation-preserving or not.
\end{prop}
\begin{proof}
By projecting $\matH^n$ orthogonally onto $l$, we give $\matH^n$ the structure of a $\matR^{n-1}$-bundle over $l$ which is preserved by $\varphi$ and hence descends to a structure of $\matR^{n-1}$-bundle over $\gamma$. The conclusion follows from the classification of vector bundles over $S^1$, see Proposition \ref{S1:bundle:prop}.
\end{proof}

A \emph{tube of radius $R$}, or a \emph{$R$-tube}, is the quotient $N_R(l)/_\Gamma$ of the closed $R$-neighbourhood $N_R(l)$ of $l$, shown in Figure \ref{tube:fig}. It is diffeomorphic to $D^{n-1} \times S^1$ or $D^{n-1}\timtil S^1$, and in particular it is compact. Note that the boundary of a tube is not totally geodesic, see Figure \ref{tube:fig}.

\begin{figure}
\begin{center}
\includegraphics[width = 4 cm] {\iftoggle{BW}{tube-BW}{tube}}
\nota{The $R$-neighbourhood $N_R(l)$ of a vertical line in the half-space model is a Euclidean cone. The boundary of the cone is not totally geodesic: the cone is however convex.}
\label{tube:fig}
\end{center}
\end{figure}

\subsection{Cusps}
We now introduce another simple type of hyperbolic manifolds. In the previous example the discrete group $\Gamma$ consisted of hyperbolic transformations fixing the same line $l$, now $\Gamma$ will consist of parabolic transformations fixing the same point at infinity.

Let $\Gamma < \Iso(\matR^{n-1})$ be a non-trivial discrete group of Euclidean isometries acting freely on $\matR^{n-1}$: the quotient $M=\matR^{n-1}/_\Gamma$ is a flat $(n-1)$-manifold. If we use the half-space model for $\matH^n$ with coordinates $(x,t)$, every element $\varphi \in \Gamma$ acts as a parabolic transformation on $\matH^n$ by sending $(x,t)$ to $(\varphi (x),t)$. The whole group $\Gamma$ is a discrete group of parabolic transformations of $\matH^n$ fixing the point $\infty$. 

The quotient $\matH^n /_\Gamma$ is naturally diffeomorphic to $M \times \matR_{>0}$. The metric tensor at the point $(x,t)$ is 
$$g_{(x,t)} = \frac{g^M_x \oplus 1}{t^2}$$
where $g^M$ is the metric tensor of the flat $M$. The manifold $\matH^n/_\Gamma$ is called a \emph{cusp}. Since $\Gamma$ contains parabolics we have $\inj M = 0$. \index{cusp}

\begin{oss}
The vertical coordinate $t$ may be parametrized more intrinsically using arc-length. As we have seen in Proposition \ref{verticale:prop}, a vertical geodesic with unit speed is parametrized as $t=e^u$. Using $u$ instead of $t$ the cusp is isometric to $M\times \matR$ with metric tensor 
$$g_{(x,u)} = (e^{-2u}g_x^M) \oplus 1.$$ 
The lengths in the flat slice $M\times u$ are shrunk or dilated by the factor $e^{-u}$.
\end{oss}

\begin{figure}
\begin{center}
\includegraphics[width = 5 cm] {\iftoggle{BW}{Pseudosphere-BW2}{Pseudosphere}}
\nota{The \emph{pseudosphere} is a surface in $\matR^3$ isometric to the union of two truncated cusps, each with constant gaussian curvature $-1$.}
\label{pseudosphere:fig}
\end{center}
\end{figure}

A \emph{truncated cusp} is a portion $N=M \times [a,+\infty)$, bounded by the Euclidean manifold $M\times a$: note that the boundary $\partial N$ is Euclidean but not totally geodesic.
The volume of a truncated cusp is strikingly simple. \index{cusp!truncated cusp}
\begin{prop} \label{truncated:volume:prop}
Let $N$ be a truncated cusp. We have
$$\Vol \big(N\big) = \frac{\Vol(\partial N)}{n-1}. $$
\end{prop}
\begin{proof}
It follows from Lemma \ref{O:lemma}.
\end{proof}

This shows in particular that a (non truncated) cusp has infinite volume.

\begin{example}
In dimension $n=2$ there is only one cusp up to isometry. The group $\Gamma<\Iso(\matR)$ is the infinite cyclic group generated by a translation $x\mapsto x+b$ and up to conjugating in $\Iso(\matH^2)$ we may take $b=1$. The cusp is diffeomorphic to $S^1\times \matR$, and the circle $S^1\times \{u\}$ has length $e^{-u}$. Some truncated cusp (but not the whole cusp!) embeds in $\matR^3$ as in Figure \ref{pseudosphere:fig}.\index{pseudosphere}
\end{example}
\begin{rem}
Pick $p\in \matH^2$. Note that a cusp and $\matH^2\setminus \{p\}$ are both hyperbolic and diffeomorphic to an open annulus $S^1\times \matR$. However, they are not isometric because the cusp is complete while $\matH^2\setminus \{p\}$ is not.
\end{rem}

We sometimes employ the word \emph{cusp} to indicate a truncated cusp, for simplicity.

\subsection{Closed geodesics}
A \emph{closed curve} in a manifold $M$ is a smooth map $\gamma\colon S^1 \to M$. A (possibly closed) curve is \emph{simple} if it is an embedding.\index{geodesic!closed geodesic}

We consider $S^1$ as the unit circle in $\matC$. A \emph{closed geodesic} in a Riemannian manifold $M$ is a smooth map $\gamma\colon S^1 \to M$ whose lift $\gamma\circ \pi\colon \matR \to M$ along the universal covering $\pi(t) = e^{i t}$ is a non-constant geodesic. 
Two closed geodesics $\gamma_1$, $\gamma_2$ that differ only by a rotation, \emph{i.e.}~such that $\gamma_1(z) = \gamma_2(ze^{it})$ for some fixed $t\in\matR$, are implicitly considered equivalent. The closed geodesics $\gamma(z)$ and $\overline\gamma(z) = \gamma(\bar z)$ are however distinct (they have opposite orientations).

\begin{prop} \label{closed:geodesic:prop}
Let $\gamma$ be a closed geodesic in a Riemannian manifold $M$. Exactly one of the following holds:
\begin{enumerate}
\item the curve $\gamma$ is simple,
\item the curve $\gamma$ self-intersects transversely in finitely many points,
\item the curve $\gamma$ wraps $k\geqslant 2$ times along a curve of type (1) or (2).
\end{enumerate}
\end{prop}
\begin{proof}
If the geodesic is not simple, it self-intersects. If it self-intersects only with distinct tangents, then (2) holds. Otherwise (3) holds.
\end{proof}

The natural number $k$ in (3) is the \emph{multiplicity} of the closed geodesic. A closed geodesic $\gamma$ of multiplicity $k$ is of type $\gamma (e^{it}) = \eta(e^{kit})$ for some geodesic $\eta$ of type (1) or (2). 

\begin{oss}
A closed geodesic on a Riemannian manifold $M$ is determined by its support, its orientation, and its multiplicity.
\end{oss}

\subsection{Closed geodesics in a hyperbolic manifold}
Closed geodesics in hyperbolic manifolds have a particularly nice behaviour. 

Let $X,Y$ be topological spaces: as usual we indicate by $[X,Y]$ the homotopy classes of continuous maps $X \to Y$. Let $X$ be path-connected. There is a natural map $\pi_1(X,x_0) \to [S^1, X]$, and the following is a standard exercise in topology.\index{$[X,Y]$} 

\begin{ex}
The map induces a bijection between the conjugacy classes in $\pi_1(X,x_0)$ and $[S^1,X]$.
\end{ex}
A simple closed curve in $X$ is \emph{homotopically trivial} if it is homotopic to a constant. As a corollary, a simple closed curve $\gamma$ is homotopically trivial if and only if it represents the trivial element in $\pi_1(X,\gamma(1))$.

On a complete hyperbolic manifold $M=\matH^n/_\Gamma$ the conjugacy classes of $\pi_1(M)$ correspond to those in $\Gamma$ and we get a correspondence
$$\big\{ {\rm conjugacy\ classes\ in\ }\Gamma\big\} \longleftrightarrow [S^1, M].$$
This correspondence works as follows: given $\varphi \in \Gamma$, pick any point $x\in\matH^n$, connect $x$ to $\gamma(x)$ with any arc, and project the arc to a closed curve in $M$.

Two conjugate elements in $\Gamma$ are of the same type (trivial, parabolic, or hyperbolic) and have the same minimum displacement. Therefore every element in $[S^1, M]$ has a well-defined type and minimum displacement.

\begin{prop} \label{unique:closed:geodesic:prop}
Let $M$ be a complete hyperbolic manifold. Every hyperbolic element of $[S^1, M]$ is represented by a unique closed geodesic, of length $d$ equal to its minimum displacement. The trivial and parabolic elements are not represented by closed geodesics.
\end{prop}
\begin{proof}
Take $M=\matH^n/_\Gamma$. A hyperbolic isometry $\varphi \in \Gamma$ has a unique invariant geodesic in $\matH^n$, its axis, which projects to a closed geodesic of length $d$. Conjugate isometries determine the same closed geodesic in $M$.

On the other hand, a closed geodesic in $M$ lifts to a segment connecting two distinct points $x_0$ and $\varphi(x_0)$ for some $\varphi\in \Gamma$ which preserves the line passing through $x_0$ and $\varphi(x_0)$: since $\varphi$ preserves a line, it is hyperbolic.
\end{proof}
We get a bijection
$$\big\{{\rm hyperbolic\ conjugacy\ classes\ in\ }\Gamma\big\} \longleftrightarrow \big\{ {\rm closed\ geodesics\ in\ } M \big\}.$$

\begin{cor} \label{unique:closed:geodesic:cor}
Let $M$ be a closed hyperbolic manifold. Every non-trivial element in $[S^1, M]$ is represented by a unique closed geodesic.
\end{cor}
\begin{proof}
Since $M$ is compact there are no parabolics.
\end{proof}

\begin{cor} \label{minimum:cor}
Let $M$ be a complete hyperbolic manifold. Every closed geodesic has the minimum length among the closed curves in its homotopy class.
\end{cor}
\begin{proof}
If $\alpha$ is a closed geodesic, its length equals the minimum displacement $d$ of a corresponding hyperbolic transformation $\varphi$. Every closed curve $\beta$ homotopic to $\alpha$ lifts to an arc connecting two points $\tilde x$ and $\varphi(\tilde x)$ that have distance at least $d$; hence $\beta$ has length at least $d$.
\end{proof}

Not only a purely topological object like a hyperbolic homotopy class of closed curves has a unique geometric nice representative, but this representative is the shortest possible one. Informally, we may think that every closed curve may be shrunk until it becomes a closed geodesic, and the negative curvature forces this closed geodesic to be unique: we will soon see that the uniqueness is lost in the elliptic and flat geometries.

It is worth recalling that a closed geodesic may not be simple. When it is simple, we now show that the closed geodesic has some nice small neighbourhoods. We defined the $R$-tubes in Section \ref{tubes:subsection}.

\begin{prop} \label{tube:prop}
The $R$-neighbourhood of a simple closed geodesic $\gamma$ in a complete hyperbolic manifold is isometric to a $R$-tube, for some $R>0$.
\end{prop}
\begin{proof}
By compactness of $\gamma$ there is a sufficiently small $R>0$ such that the $R$-neighbourhood of $\gamma$ lifts to disjoint $R$-neighbourhoods of its geodesic lifts in $\matH^n$. Hence their quotient is a $R$-tube.
\end{proof}


\section{The Margulis Lemma}
We state and prove the \emph{Margulis Lemma}, that concerns arbitrary discrete groups in Lie groups and more specifically in $\Iso(\matH^n)$. The lemma implies that there is a constant $\varepsilon>0$, which depends only on the dimension $n$, such that the $\varepsilon$-thin part of any complete hyperbolic $n$-manifold consists of truncated cusps and tubes only. The $\varepsilon$-thin part is by definition the set of all the points with injectivity radius smaller than $\varepsilon$.

The proof of the lemma for general Lie groups is surprisingly simple and elegant; its application to the hyperbolic case needs however a more technical argument. We start by exposing some preliminary facts on the isometries of the hyperbolic space that are of independent interest.

\subsection{Isometries that commute or generate discrete groups} \label{that:commute:subsection}
Two isometries of $\matH^n$ that commute or generate a discrete group must be of a particular kind. We let $\Fix(\varphi)$ denote the fixed points in $\overline{\matH^n}$ of an isometry $\varphi$. 

\begin{lemma} \label{commute:lemma}
Let $\varphi_1,\varphi_2\in\Iso(\matH^n)$ be two hyperbolic or parabolic isometries. If they commute then $\Fix(\varphi_1) = \Fix(\varphi_2)$.
\end{lemma}
\begin{proof}
If they commute, the map $\varphi_1$ acts on $\Fix(\varphi_2)$ and vice versa. If $\varphi_2$ is hyperbolic, then $\Fix(\varphi_2)=\{p,q\}$ and $\varphi_1$ fixes the line with endpoints $p$ and $q$, hence is again hyperbolic with $\Fix(\varphi_1)=\{p,q\}$. If $\varphi_1$ and $\varphi_2$ are parabolic then they have the same fixed point $\Fix(\varphi_1) = \Fix(\varphi_2)$.
\end{proof}

\begin{lemma} \label{fissi:lemma}
Let $\varphi_1,\varphi_2\in\Gamma$ be two non-trivial isometries in a discrete group $\Gamma < \Iso (\matH^n)$ that acts freely on $\matH^n$. Either $\Fix(\varphi_1)\cap \Fix(\varphi_2) = \emptyset$, or one of the following holds:
\begin{itemize}
\item $\varphi_1$ and $\varphi_2$ are parabolics with the same fixed point,
\item $\varphi_1$ and $\varphi_2$ are powers of the same hyperbolic $\varphi \in \Gamma$.
\end{itemize}
\end{lemma}
\begin{proof}
Suppose that $\varphi_1$ and $\varphi_2$ have some common fixed points. We first consider the case where $\varphi_1$ is hyperbolic and $\varphi_2$ is parabolic. We pick the half-space model and suppose $\Fix(\varphi_1) = \{0,\infty\}$ and $\Fix(\varphi_2) = \{\infty\}$. 

Proposition \ref{tre:prop} says that 
$$\varphi_1(x,t) = \lambda(Ax,t), \quad \varphi_2(x,t) = (A'x+b, t)$$ 
with $A, A' \in \On(n-1)$ and $\lambda \neq 1$. Hence
\begin{align*}
\varphi_1^n\circ \varphi_2\circ  \varphi_1^{-n}(x,t) & = \varphi_1^n\big(A'(\lambda^{-n}A^{-n}x)+b,\lambda^{-n}t) \\
& = \big(A^nA'A^{-n}x + \lambda^n A^nb , t\big).
\end{align*}
Up to interchanging $\varphi_1$ and $\varphi_1^{-1}$ we may suppose $\lambda < 1$ and get
$$\lim_{n\to \infty} \varphi_1^n \circ \varphi_2 \circ \varphi_1^{-n}(0,t) = 
\lim_{n\to \infty} \left(\lambda^nA^nb, t\right) = (0,t).$$
A contradiction since $\Gamma$ is discrete.

We now suppose that both $\varphi_1$ and $\varphi_2$ are hyperbolic, with $\Fix(\varphi_1) = \{a,\infty\}$ and $\Fix (\varphi_2) = \{b,\infty\}$. The isometries $\varphi_1$ and $\varphi_2$ permute the horospheres centred at $\infty$ and
$$[\varphi_1,\varphi_2] = \varphi_1\circ \varphi_2 \circ \varphi_1^{-1} \circ \varphi_2^{-1} \in \Gamma$$
fixes every horizontal horosphere. Hence the commutator is parabolic or trivial: the first case is excluded by the previous discussion, in the second case we have $a=b$ by Lemma \ref{commute:lemma}. 

Both $\varphi_1$ and $\varphi_2$ have the same axis $l$, and since they generate a discrete group $\Gamma_0 < \Gamma$ they are both powers of some hyperbolic $\varphi\in\Gamma_0$ with that axis. To prove this, note that $\Gamma_0$ acts effectively on $l$ as a discrete group of translations, hence $\Gamma_0 \isom \matZ$. \end{proof}

Two isometries are contained in some discrete group if and only if they generate a discrete group, so the previous lemma is actually a fact on pairs of isometries that generate discrete groups.
The two lemmas have important geometric consequences.

\begin{cor}  \label{axis:cor}
Let $\matH^n/_\Gamma$ be a complete hyperbolic manifold. The axis in $\matH^n$ of two hyperbolic isometries in $\Gamma$ are either incident or ultra-parallel (not asymptotically parallel). 
\end{cor}

\begin{cor} \label{ZxZ:cor}
Let $\matH^n/_\Gamma$ be a complete hyperbolic manifold. Every subgroup of $\Gamma$ isomorphic to $\matZ\times\matZ$ consists of parabolic elements fixing the same point at infinity.
\end{cor}

\begin{cor} \label{no:ZxZ:cor}
The fundamental group $\pi_1(M)$ of a closed hyperbolic manifold $M$ does not contain subgroups isomorphic to $\matZ\times \matZ$.
\end{cor}

\begin{cor} \label{torus:no:hyperbolic:cor}
The torus has no hyperbolic structure.
\end{cor}

\subsection{The Margulis Lemma}
We now introduce the main character of this chapter, the Margulis Lemma. We first state a general version for Lie groups and then turn more specifically to $\Iso(\matH^n)$.

\begin{lemma} \label{G:lemma}
Let $G$ be a Lie group. There is a neighbourhood $U$ of $e\in G$ such that every discrete subgroup $\Gamma < G$ generated by some elements in $U$ is nilpotent.
\end{lemma}
Note that the discreteness of $\Gamma$ is essential here, since every connected Lie group $G$ is generated by arbitrarily small elements, and $G$ needs not to be nilpotent.

\begin{proof}
Consider the commutator map 
$$[\ ,\ ]\colon G\times G \to G$$ 
that sends $(g,h)$ to the commutator $[g,h]$. The map is smooth and sends $G\times \{e\}$ and $\{e\}\times G$ to the point $e$. Therefore its differential at $(e,e)$ vanishes. We identify a neighbourhood of $e$ in $G$ with $\matR^n$, so that the commutator map near $(e,e)$ can be read as 
$$[\ ,\ ]\colon U \times U \longrightarrow \matR^n$$
for some neighbourhood $U\subset \matR^n$ of the origin.
Since its differential at $(0,0)$ vanishes, up to restricting $U$ we may suppose that the map is $\frac 12$-Lipschitz. Therefore for all $x,y \in U$ we get
$$\|[x,y]\| < \|[x,0]\| + \frac 12\|y\| = \frac 12\|y\|, \quad \|[x,y]\| < \|[0,y]\| + \frac 12\|x\| = \frac 12\|x\|$$
and hence
$$\|[x,y]\| < \frac 12 \min \big\{\|x\|, \|y\|\big\}.$$
This implies that for every smaller neighbourhood $V\subset U$ of $0$ there is a $k>0$ such that
$$\underbrace{[U, [U, \ldots [U,U]]\cdots ]}_k \subset V.$$
We now turn back to $G$. Let $\Gamma$ be discrete and generated by some elements $S\subset U$. We choose a smaller neighbourhood $V$ of $e$ such that $V \cap \Gamma = \{e\}$, hence for any $a_1,\ldots, a_k,b\in S$ we get $[a_1,[a_2, \ldots[a_k,b]]\cdots ] \in V\cap\Gamma =\{e\}$. Proposition \ref{criterion:nilpotent:prop} says that $\Gamma$ is nilpotent.
\end{proof}

We now want to refine this lemma when $G=\Iso(\matH^n)$. Let $P$ be a property of groups, like being abelian, nilpotent, etc. A group is \emph{virtually $P$} if it has a finite-index subgroup which is $P$.\index{Margulis Lemma} 

\begin{lemma}[Margulis Lemma] \label{Margulis:lemma}
In every dimension $n\geqslant 2$ there is a constant $\varepsilon_n>0$ such that for all $x\in \matH^n$, every discrete group $\Gamma<\Iso(\matH^n)$ generated by elements that move $x$ at distance smaller than $\varepsilon_n$ is virtually nilpotent.
\end{lemma}
\begin{proof}
It suffices to prove the theorem for a fixed $x\in \matH^n$ since the isometries of $\matH^n$ act transitively on points. By Lemma \ref{G:lemma} there is a neighbourhood $U$ of $e$ in $G=\Iso(\matH^n)$ such that every discrete group generated by some elements in $U$ is nilpotent. 

\begin{figure}
\begin{center}
\includegraphics[width = 5 cm] {\iftoggle{BW}{Margulis-BW}{Margulis}}
\nota{The Margulis Lemma for Lie groups furnishes a nice neighbourhood $U$ of $e\in\Iso(\matH^n)$, but here we need a nice neighbourhood $W$ of the whole compact stabiliser $\On(n)$, drawn as a \iftoggle{BW}{dark grey}{blue} line.}
\label{Margulis:fig}
\end{center}
\end{figure}

Let $V_\varepsilon\subset\Iso(\matH^n)$ be the set of all the isometries that move $x$ at distance smaller than $\varepsilon>0$. It is relatively compact. As $\varepsilon>0$ varies, the sets $V_\varepsilon$ form a neighbourhood system for the compact stabiliser $G_x <\Iso(\matH^n)$ of $x$, isomorphic to $\On(n)$. Pick a $V_\varepsilon$. Being relatively compact, it can be covered by some $m-1$ translates of $U$. See Figure \ref{Margulis:fig}.

Let now $W=W^{-1}$ be a symmetric neighbourhood of $G_x$ satisfying $W^m \subset V_\varepsilon$. Let $\Gamma$ be a discrete group generated by some elements $S\subset W$. We want to prove that the nilpotent subgroup $\Gamma_U$ generated by $\Gamma \cap U$ has index at most $m$ in $\Gamma$. This concludes the proof, for it suffices to take a $\varepsilon_n$ such that $V_{\varepsilon_n} \subset W$.

Suppose by contradiction that $\Gamma/_{\Gamma_{U}}$ contains more than $m$ cosets. Let $r(i)$ be the number of cosets represented by elements in $\Gamma$ that are products of at most $i$ generators in $S$. If $r(i)=r(i+1)$ for some $i$, then it is easy to deduce that $r(i)$ stabilises forever. This easily implies that $r(m)\geqslant m$. Therefore there are $m$ elements in $W^m$ belonging to $m$ distinct cosets. By hypothesis two of them, say $g$ and $h$, belong to the same translate of $U$. Therefore $gh^{-1} \in \Gamma\cap U$, contradicting the fact that $g$ and $h$ lie in distinct cosets.
\end{proof}

For a discrete group $\Gamma < \Iso(\matH^n)$ and a point $x\in \matH^n$, we denote by $\Gamma_\varepsilon (x)< \Gamma$ the subgroup generated by all elements $g\in \Gamma$ that move $x$ of a distance smaller than $\varepsilon$. We have proved that $\Gamma_{\varepsilon_n} (x)$ is virtually nilpotent for every $x\in\matH^n$ and every discrete $\Gamma$.

A \emph{Margulis constant} in a fixed dimension $n$ is any number $\varepsilon_n>0$ for which Lemma \ref{Margulis:lemma} holds. Every sufficiently small number is a Margulis constant.\index{Margulis constant} 

\subsection{Elementary groups}
A non-trivial discrete group $\Gamma < \Iso(\matH^n)$ is \emph{elementary} if it preserves a finite set of points in $\overline{\matH^n}$.\index{group!elementary group}

\begin{prop} \label{elementary:prop}
An elementary $\Gamma$ acting freely on $\matH^n$ is:
\begin{itemize}
\item generated by a hyperbolic isometry, or
\item generated by parabolic isometries having the same fixed point at $\infty$.
\end{itemize}
\end{prop}
\begin{proof}
Since $\Gamma$ contains no elliptics, every non-trivial element fixes one or two points at infinity in $\overline{\matH^n}$, and no other finite set of points (because of Proposition \ref{iterate:prop}). We conclude using Lemma \ref{fissi:lemma}.
\end{proof}

\begin{prop} \label{elementary:index:prop}
Let $\Gamma<\Iso(\matH^n)$ be a discrete group acting freely on $\matH^n$. If $\Gamma'<\Gamma$ has finite index and is elementary, then $\Gamma$ also is.
\end{prop}
\begin{proof}
We know that every element in $\Gamma'$ is either hyperbolic with axis $l$, or parabolic with fixed point $p\in\partial\matH^n$. If $\varphi \in \Gamma$ then $\varphi^k\in \Gamma'$ for some $k$: hence $\varphi$ is also of that type.
\end{proof}

\begin{cor}
Let $\Gamma<\Iso(\matH^n)$ be discrete and acting freely. If $\Gamma$ is virtually nilpotent, it is either trivial or elementary.
\end{cor}
\begin{proof}
We know that $\Gamma$ contains a finite-index nilpotent subgroup $H$. If $H$ is trivial then $\Gamma$ is finite and hence trivial. If $H$ is non-trivial, it has a non-trivial centre by Proposition \ref{center:prop}. Lemma \ref{commute:lemma} then implies that all the elements in $H$ have the same fixed points and so $H$ is elementary. Therefore $\Gamma$ is elementary by Proposition \ref{elementary:index:prop}
\end{proof}

Let $\varepsilon_n$ be a Margulis constant. We can strengthen the Margulis Lemma.

\begin{cor}
Let $\Gamma<\Iso(\matH^n)$ be discrete and acting freely. For every point $x\in \matH^n$ the subgroup $\Gamma_{\varepsilon_n}(x)$ is either trivial or elementary.
\end{cor}

\subsection{Thick-thin decomposition}
A notable geometric consequence of the Margulis Lemma is that every complete hyperbolic manifold decomposes into a thick and a thin part, where the thin part consists only of particularly simple star-shaped pieces. This decomposition is known as the \emph{thick-thin decomposition} of hyperbolic manifolds.\index{thick-thin decomposition}

We define a \emph{star-shaped set} centred at $p\in\partial\matH^n$ to be any subset $U\subset\matH^n$ that intersects every half-line pointing to $p$ in a half-line. For instance, a horoball is star-shaped. A \emph{star-shaped neighbourhood} of a line $l\subset\matH^n$ is any neighbourhood $V$ of $l$ that intersects every line orthogonal to $l$ into a connected set. For instance, a $R$-neighbourhood of $l$ is star-shaped.

These definitions pass to quotients. A \emph{star-shaped cusp neighbourhood} is the quotient $U/_\Gamma$ of a $\Gamma$-invariant star-shaped set $U$ centred at $p\in\partial \matH^n$ via a discrete group $\Gamma$ of parabolic transformations fixing $p$ and acting freely. Analogously, a \emph{star-shaped simple closed geodesic neighbourhood} is the quotient $V/_\Gamma$ of a $\Gamma$-invariant star-shaped neighbourhood $V$ of $l$ via a discrete group $\Gamma\isom \matZ$ of hyperbolic transformations with axis $l$.

The truncated cusps and $R$-tubes studied in Section \ref{tubes:cusps:section} are particularly nice star-shaped cusp and geodesic neighbourhoods. 

Let $\varepsilon_n>0$ be a Margulis constant. We define $M_{[\varepsilon_n,\infty)}$ and $M_{(0,\varepsilon_n]}$ respectively as the set of all points $x \in M$ having $\inj_x M \geqslant \frac{\varepsilon_n}2$, and as the closure of the complementary set $M\setminus M_{[\varepsilon_n,\infty)}$. They form respectively the \emph{thick} and \emph{thin part} of $M$. 

\begin{oss} \label{exclude:thin:oss}
We do not define the thin part simply as the set of all points $x$ having $\inj_x M \leqslant \frac{\varepsilon_n}2$ because we want to discard the degenerate and unlucky case of a closed geodesic $\gamma$ having length precisely $\varepsilon_n$. The injectivity radius would be $\frac{\varepsilon_n}2$ at the points in $\gamma$ and strictly bigger than $\frac{\varepsilon_n}2$ near $\gamma$. With our definition, the geodesic $\gamma $ is contained in $M_{[\varepsilon_n, +\infty)}$ and not in $M_{(0,\varepsilon_n]}$.
\end{oss}

The following theorem is arguably the most important structural result on complete hyperbolic manifolds of any dimension $n$.

\begin{teo}[Thick-thin decomposition] Let $M$ be a complete hyperbolic $n$-manifold. The thin part $M_{(0,\varepsilon_n]}$ consists of a disjoint union of star-shaped neighbourhoods of cusps and of simple closed geodesics of length $< \varepsilon_n$.
\end{teo}
\begin{proof}
We have $M=\matH^n/_\Gamma$. For every isometry $\varphi\in \Gamma $ we define
$$S_\varphi(\varepsilon) = \big\{ x\in\matH^n\ \big|\ d(\varphi(x),x)\leqslant \varepsilon\big\}\subset \matH^n.$$
By Proposition \ref{d:prop} the thin part $M_{(0,\varepsilon_n]}$ is the image of the set 
\begin{align*}
 S & = \big\{x\in \matH^n\ \big|\ \exists \varphi\in \Gamma, \varphi\neq \id {\rm\ such\ that\ } d(\varphi(x),x)\leqslant \varepsilon_n\big\} \\
 & = \bigcup_{\varphi\in \Gamma, \varphi \neq \id} S_\varphi(\varepsilon_n).
 \end{align*}
More precisely, we should exclude the hyperbolic $\varphi \in \Gamma$ with $d(\varphi)=\varepsilon_n$, see Remark \ref{exclude:thin:oss}, but this is not an important point and we ignore it. It is easy to check that $S_\varphi(\varepsilon)$ is star-shaped, centred at a $p\in \partial \matH^n$ or at a line $l$ according to whether $\varphi$ is parabolic fixing $p$ or hyperbolic fixing $l$.

Suppose that $x\in S_\varphi(\varepsilon_n) \cap S_\psi(\varepsilon_n)$ for some non-trivial isometries $\varphi, \psi \in \Gamma$. By the Margulis Lemma both $\varphi$ and $\psi$ belong to the elementary group $\Gamma_{\varepsilon_n}(x)$ and hence by Proposition \ref{elementary:prop} both $\varphi$ and $\psi$ are either parabolic fixing the same point $p$ at infinity or hyperbolic fixing the same line $l$.

Therefore every connected component $S_0$ of $S$ is the union of all $S_\varphi(\varepsilon_n)$ where $\varphi$ varies in some maximal elementary subgroup $\Gamma_0<\Gamma$ of parabolics fixing the same point $p$ or hyperbolics fixing the same line $l$. The set $S_0$ is a union of star-shaped sets centred at $p$ or $l$ and is hence also star-shaped. 

The group $\Gamma$ preserves $S$ and the only isometries in $\Gamma$ that preserve $S_0$ are those in $\Gamma_0$, therefore the quotient $M_{(0,\varepsilon_n]} = S/_\Gamma$ consists of star-shaped neighbourhoods of cusps and of simple closed geodesics.
\end{proof}

Star-shaped neighbourhoods are particularly nice in low dimensions.

\begin{prop}
Let $M$ be a complete orientable hyperbolic manifold of dimension $n\leqslant 3$. The thin part $M_{(0,\varepsilon_n]}$ consists of truncated cusps and $R$-tubes.
\end{prop}
\begin{proof}
Pick a non-trivial $\varphi \in \Iso^+(\matH^n)$. We check that $S_\varphi(\varepsilon)$ is either empty, or a $R$-neighbourhood of a line $l$, or a horoball, for all $\varepsilon>0$. This proves the proposition.

Suppose that $\varphi$ is hyperbolic with axis $l$. The distance $d(x,\varphi(x))$ depends only on $d(x,l)$ since all the orientation-preserving hyperbolic transformations with axis $l$ commute with $\varphi$ and act transitively on the points at fixed distance from $l$. 
It is easy to check that $d(x,\varphi(x))$ increases with $d(x,l)$ and hence $S_\varphi(\varepsilon)$ is either empty or a $R$-neighbourhood of $l$.

If $\varphi$ is parabolic, it acts on each horosphere centred at $p$ like a
Euclidean fixed-point-free orientation-preserving isometry on $\matR^{n-1}$: this must be a translation when $n\leqslant 3$. Therefore $d(x,\varphi(x))$ depends only on the horosphere $O$ containing $x$ and decreases as $O$ moves towards $p$. Therefore $S_\varphi(\varepsilon)$ is a horoball.
\end{proof}

\begin{figure}
\begin{center}
\includegraphics[width = 11.5 cm] {\iftoggle{BW}{thick_thin-BW}{thick_thin}}
\nota{The thick-thin decomposition of a complete hyperbolic surface: the thin part (\iftoggle{BW}{light grey}{yellow}) consists of truncated cusps and neighbourhoods of short geodesics\iftoggle{BW}{}{ (blue)}.}
\label{thick_thin:fig}
\end{center}
\end{figure}

See the picture in Figure \ref{thick_thin:fig}. 

\begin{cor} \label{simple:small:geodesics:cor}
Let $M$ be a complete hyperbolic $n$-manifold. The closed geodesics in $M$ of length $< \varepsilon_n$ are simple and disjoint.
\end{cor}
\begin{proof}
These closed geodesics lie in the thin part. Star-shaped cusp neighbourhoods contain no closed geodesics, and each star-shaped geodesic neighbourhood contains only one closed geodesic, its core, which is simple.
\end{proof}

\subsection{Finite-volume hyperbolic manifolds} \label{finite-volume:subsection}
A specific feature of hyperbolic geometry is the existence of complete hyperbolic manifolds that have finite volume without being compact. The thin-thick decomposition furnishes a nice topological description of such manifolds.

\begin{prop}
A complete hyperbolic manifold $M$ has finite volume if and only if its thick part is compact.
\end{prop}
\begin{proof}
If the thick part is not compact, it contains an infinite number of points that stay pairwise at distance greater than $\varepsilon_n$. The open balls of radius $\frac{\varepsilon_n}2$ centred at these points are embedded and disjoint and all have the same volume: therefore their union has infinite volume.

If the thick part is compact, it has finite volume. Its boundary is also compact, and hence has finitely many connected components. Therefore the thin part
consists of finitely many star-shaped neighbourhoods of cusps and closed geodesics, each with compact boundary. Each such object has finite volume (because it is contained in a bigger abstract truncated cusp with compact base, or a $R$-tube, which has finite volume).
\end{proof}

\begin{cor} \label{tame:cor}
Every complete finite-volume hyperbolic manifold $M$ is diffeomorphic to the interior of a compact manifold $N$ with boundary. The boundary $\partial N$ consists of manifolds that admit some flat structure.
\end{cor}
\begin{proof}
We have seen in the previous proof that the thick part of $M$ is compact, and the thin part consists of finitely many star-shaped neighbourhoods of cusps and closed geodesics, each with compact boundary.

Every cusp neighbourhood has compact base and hence contains a smaller truncated cusp, diffeomorphic to $X\times [0,1)$ for some closed flat manifold $X$. The complement in $M$ of these truncated cusps is compact: it is obtained from the thick part by adding finitely many neighbourhoods of closed geodesics (that are compact) and compact portions of cusp neighbourhoods. Each truncated cusp is diffeomorphic to $X\times [0,1)$ and can be compactified by adding $X \times 1$. The resulting manifold $N$ is compact with boundary.
\end{proof}

Every boundary component $X$ of $N$ inherits a flat structure, uniquely determined up to rescaling (different truncations modify $X$ only by rescaling.) As already mentioned, a truncated cusp is often called simply a \emph{cusp}: for instance we say that the surface sketched in Figure \ref{thick_thin:fig} has two cusps.

We briefly discuss the effects of Corollary \ref{tame:cor} on low-dimensional manifolds. 
Let $M$ be a complete hyperbolic manifold of finite volume of dimension $n$. If $n=2$, the manifold $M$ is diffeomorphic to the interior of a compact surface $N$ with boundary. Every boundary component of $N$ is of course a circle.

If $n=3$, the manifold $M$ is diffeomorphic to the interior of a compact 3-manifold $N$ with boundary, and every boundary component $X$ of $N$ is a flat surface: if $M$ is orientable, then $\partial N$ also is, and by Proposition \ref{flat:is:torus:prop} it consists of tori.

\subsection{Geodesic boundary and cusps}
Most of the arguments of this chapter extend to hyperbolic manifolds with compact geodesic boundary.

\begin{prop} \label{geodesic:M:N:prop}
Every complete finite-volume hyperbolic manifold $M$ with compact geodesic boundary is diffeomorphic to a compact manifold $N$ with some boundary components removed. The removed components have a flat structure. The remaining components form the geodesic boundary of $M$ and hence have a hyperbolic structure.
\end{prop}
\begin{proof}
Double $M$ along the geodesic boundary to obtain a finite-volume hyperbolic manifold to which Corollary \ref{tame:cor} applies. 
\end{proof}

We briefly discuss the effects in low dimensions. In dimension $n=2$ the boundary $\partial N$ consists of circles: some are geodesic components of $\partial M$, while some others are removed and correspond to cusps. In dimension $n=3$, if $M$ is orientable the removed boundary of $\partial N$ consists of tori (the cusps), while the geodesic boundary of $M$ cannot contain any torus: a torus has no hyperbolic structure by Corollary \ref{torus:no:hyperbolic:cor}.

We add some information on fundamental groups. Let the manifolds $M$ and $N$ be as in Proposition \ref{geodesic:M:N:prop}.

\begin{prop} \label{pi1:M:prop}
For every boundary component $X$ of $N$ the homomorphism $\pi_1(X)\to\pi_1(N)$ induced by inclusion is injective. Two homotopically non-trivial closed curves in distinct boundary components of $N$ are not freely homotopic in $N$.
\end{prop}
\begin{proof}
We have $M=C/_\Gamma$ for some convex $C\subset \matH^3$ with boundary consisting of hyperplanes. If $X$ is a (cusp) flat component of $\partial N$, it is isometric to $\matR^{n-1}/_{\Gamma_X}$ and the map $\Gamma_X=\pi_1(X) \to \pi_1(N)=\Gamma$ is injective since it sends a Euclidean isometry to a corresponding parabolic isometry.
If $X$ is geodesic, it lifts to a hyperplane in $\partial C$, which is simply connected: hence $\pi_1(X) \to \pi_1(N)$ is injective.

A homotopically non-trivial closed curve $\gamma\subset X$ determines a conjugacy class of hyperbolic or parabolic transformations in $\Gamma$. Every element in this class preserves a unique boundary hyperplane or a point at infinity of $C$ that projects back to $X$, so the conjugacy class determines $X$. Therefore distinct boundary components contain different conjugacy classes. 
\end{proof}

We note that a finite-volume hyperbolic manifold $M$ may have non-compact geodesic boundary: an ideal polygon in $\matH^2$ is a simple example. 

\section{Geodesic spectrum, isometry groups, and finite covers} \label{geodesic:spectrum:section}
We now study the geodesic spectrum and the isometry group of hyperbolic manifolds. The geodesic spectrum of $M$ is the set of the lengths of all the closed geodesics in $M$. We prove that a finite-volume $M$ has a discrete geodesic spectrum (with finite multiplicities) and a finite isometry group.\index{geodesic spectrum}

Then we turn to finite covers and derive some geometric consequences from the residually finiteness of fundamental groups.

We are almost exclusively interested in finite-volume hyperbolic manifolds. The general strategy when proving something about finite-volume hyperbolic manifolds is the following: we first suppose that the manifold is compact for simplicity, and then we adapt the proof to the non-compact case by looking at what happens to the cusps.

\subsection{Geodesic spectrum} \label{geodesic:spectrum:subsection}
We start by proving the following.
\begin{prop} Let $M$ be a finite-volume complete hyperbolic manifold. For every $L>0$ there are finitely many closed geodesics in $M$ shorter than $L$.
\end{prop}
\begin{proof}
Suppose that there are infinitely many closed geodesics shorter than $L$. We know that $M$ decomposes into a compact part and a finite union of truncated cusps. The compact part has finite diameter $D$. 

Every closed geodesic intersects the compact part because a cusp contains no closed geodesic. Therefore we can fix a basepoint $x_0\in M$ and connect $x_0$ to these infinitely many geodesics of length $<L$ with arcs shorter than $D$. We use these arcs to freely homotope the geodesics into loops based at $x_0$ of length $<L+2D$, and lift the loops to arcs in $\matH^n$ starting from some basepoint $\tilde{x}_0\in \matH^n$.

If two such arcs end at the same point, the corresponding closed geodesics in $M$ are freely homotopic: this is excluded by Proposition \ref{unique:closed:geodesic:prop}, hence these endpoints are all distinct. The orbit of $\tilde{x}_0$ now contains infinitely many points in the ball $B(\tilde{x}_0, L+2D)$, a contradiction because the orbit is discrete.
\end{proof}

The lengths of the closed geodesics in $M$ form a discrete subset of $\matR$ called the \emph{geodesic spectrum} of $M$. Let $\ell_1(M)>0$ be the minimum of the spectrum of $M$: a \emph{shortest closed geodesic} in $M$ is a closed geodesic of shortest length $\ell_1(M)$. It is not necessarily unique (not even if considered up to orientation reversal), but there are finitely many of them.\index{shortest geodesic}

\begin{prop}
If $M$ is a closed hyperbolic manifold, then 
$$\inj M = \frac 12 \cdot \ell_1(M)$$
and every shortest closed geodesic is simple.
\end{prop}
\begin{proof}
The length of a closed geodesic is the minimum displacement of the corresponding hyperbolic transformation, so Corollary \ref{inj:cor} gives the equality.

Let now $\gamma$ be a shortest closed geodesic. Consider Proposition \ref{closed:geodesic:prop}:
if $\gamma$ wraps multiple times along a curve $\eta$, then $\eta$ is shorter than $\gamma$, a contradiction. If $\gamma$ self-intersects transversely at some point $p$, we can split $\gamma$ naturally as $\gamma_1*\gamma_2$ in $\pi_1(M,p)$, where both $\gamma_1$ and $\gamma_2$ are shorter than $\gamma$. Either $\gamma_1$ or $\gamma_2$ is non-trivial in $\pi_1(M,p)$ and is hence hyperbolic, but it has length smaller than $\ell_1(M)$: this contradicts Corollary \ref{minimum:cor}.
\end{proof}

\begin{rem}
A shortest geodesic may be non-simple when $M$ has cusps! By doubling an ideal triangle along its boundary we construct a hyperbolic surface called the \emph{thrice-punctured sphere} which has three cusps and contains various closed geodesics, none of which is simple. See Chapter \ref{surfaces:chapter}.
\end{rem}

\subsection{Isometry group}
We now study the isometry group $\Iso(M)$ of a hyperbolic manifold $M$.
Recall that the \emph{normaliser} $N(H)$ of a subgroup $H < G$ is the set of elements $g\in G$ such that $gH=Hg$. It is the biggest subgroup of $G$ containing $H$ as a normal subgroup.
The isometry group $\Iso(M)$ has an algebraic representation.\index{normaliser}

\begin{prop}
Let $M=\matH^n/_\Gamma$ be a complete hyperbolic manifold. There is a natural isomorphism 
$$\Iso(M) \Isom N(\Gamma)/_\Gamma.$$
\end{prop}
\begin{proof}
Every isometry $\varphi\colon M \to M$ lifts to an isometry $\tilde \varphi$ 
$$
\xymatrix{ 
\matH^n\ar@{.>}[r]^{\tilde \varphi} \ar[d]_\pi & \matH^n \ar[d]^\pi \\
M\ar[r]_\varphi & M
}
$$
such that $\tilde\varphi \Gamma = \Gamma \tilde\varphi$: hence $\tilde\varphi \in N(\Gamma)$. The lift is uniquely determined up to left- or right-multiplication by elements in $\Gamma$, hence we get a homomorphism
$$\Iso(M) \to N(\Gamma)/_\Gamma$$
which is clearly surjective (every element in $N(\Gamma)$ determines an isometry) and injective (if $\tilde\varphi\in\Gamma$ then $\varphi=\id$).
\end{proof}

Recall that the \emph{centraliser} of $H< G$ is the set of elements $g \in G$ such that $gh=hg$ for all $h$. It is a subgroup of $G$.\index{centraliser}

\begin{ex} \label{centralizer:ex}
Let $M=\matH^n/_\Gamma$ be a finite-volume hyperbolic manifold. The centralizer of $\Gamma$ in $\Iso(\matH^n)$ is trivial.
\end{ex}

\subsection{Outer automorphism group}
The \emph{automorphism group} of a group $G$ is the group $\Aut(G)$ of all the isomorphisms $G\to G$. The \emph{inner automorphisms} are those isomorphisms of type $g\mapsto hgh^{-1}$ for some $h\in G$, and they form a normal subgroup $\Int(G)\triangleleft \Aut(G)$. The quotient
$$\Out(G) = \Aut(G)/_{\Int(G)}$$
is called the \emph{outer automorphism group} of $G$.\index{group!automorphism group}\index{group!outer automorphism group}

If $x_0, x_1$ are two points in a path-connected topological space $X$ there is a non-canonical isomorphism
$\pi_1(X,x_0) \to \pi_1(X,x_1)$, unique only up to post-composing with an inner automorphism. Therefore there is a \emph{canonical} isomorphism $\Out(\pi_1(X,x_0))\to\Out(\pi_1(X,x_1))$. Hence $\Out(\pi_1(X))$ depends very mildly on the basepoint.

The group $\Omeo(X)$ of all homeomorphisms of $X$ does not act directly on $\pi_1(X)$ because of the inner-automorphism ambiguity, but we get a natural homomorphism 
$$\Omeo (X) \longrightarrow \Out(\pi_1(X))$$
which is neither injective nor surjective in general. We note that homotopic self-homeomorphisms give rise to the same element in $\Out(\pi_1(X))$.

\subsection{Finite isometry groups}
We turn back to hyperbolic manifolds. 

\begin{prop} \label{injective:prop}
If $M$ is a finite-volume complete hyperbolic manifold, the natural map 
$$\Iso(M) \longrightarrow \Out(\pi_1(M))$$
is injective.
\end{prop}
\begin{proof}
Set $M=\matH^n/_\Gamma$, identify $\Gamma$ with $\pi_1(M)$ and $\Iso(M)$ with $N(\Gamma)/_\Gamma$. With these identifications the map
$$N(\Gamma)/_\Gamma \longrightarrow \Out(\Gamma)$$
is just the conjugacy action that sends $h\in N(\Gamma)$ to the automorphism $g \mapsto h^{-1}gh$ of $\Gamma$. This is an inner automorphism if and only if there is a $f \in\Gamma$ such that $h^{-1}gh = f^{-1}gf$ for all $g\in\Gamma$, that is if $hf^{-1}$ commutes with $g$ for all $g\in\Gamma$. Exercise \ref{centralizer:ex} gives $h=f\in\Gamma$ and hence the map is injective.
\end{proof}

\begin{cor} \label{isometries:not:homotopic:cor}
Let $M$ be a finite-volume complete hyperbolic manifold. Distinct isometries of $M$ are not homotopic.
\end{cor}

This is a quite strong fact, from which we deduce the following.

\begin{cor} \label{finite:isometry:group:cor}
The isometry group of every finite-volume complete hyperbolic manifold $M$ is finite. 
\end{cor}
\begin{proof}
We know that $M$ is diffeomorphic to the interior of some compact $N$ having $k\geqslant 0$ boundary components. Let $C_1,\ldots, C_k \subset M$ be disjoint truncated cusps, all of the same small volume $V>0$. Every isometry of $M$ permutes these truncated cusps and fixes their complement $C = M\setminus \interior {C_1\cup \ldots \cup C_k}$.

The set $C$ is compact and preserved by the Lie group $\Iso(M)$, which is hence compact. To show that it is finite it suffices to prove that it is discrete. Suppose that a sequence of isometries $\varphi_i$ converges to the identity: we will deduce that $\varphi_i$ is homotopic to $\id$ for sufficiently large values of $i$, a contradiction.

For any $\varepsilon>0$ there is a $\varphi_i$ that moves all points of $C$ at distance $<\varepsilon$. 
If $\varepsilon>0$ is sufficiently small, the isometry $\varphi_i$ preserves each $C_i$, 
it restricts to an isometry of the flat torus $X_i = \partial C_i$ and its extension to $C_i$ is just a smaller rescaling of $\varphi_i|_{X_i}$ at every flat torus leaf in the cusp. 

Let also $\varepsilon>0$ be sufficiently small, so that $\varepsilon <\inj_x M$ for every $x\in C$. We get $d(x,\varphi_i(x))<\inj_xM$ for all $x\in C$, and since $\varphi_i$ acts on each $C_i$ just by rescaling its action on $X_i$, we get the same inequality for every $x\in M$.

For every $x\in M$ there is a unique geodesic $\gamma_x$ of length $d(x,\varphi_i(x))$ connecting $x$ and $\varphi_i(x)$. The geodesics $\gamma_x$ vary continuously with $x$ and we can use them to define a homotopy between $\varphi_i$ and $\id$: a contradiction.
\end{proof}

The finiteness of the isometry groups is peculiar to the hyperbolic world: we will see in Remark \ref{torus:isometries:oss} that the isometry group of a flat torus is infinite.

\subsection{Finite covers} \label{finite:covers:subsection}
Let $M=\matH^n/_\Gamma$ be a complete finite-volume hyperbolic manifold. How many finite covers $\tilde M \to M$ are there above $M$? Quite a lot, thanks to the following.

\begin{prop} \label{RF:hyperbolic:prop}
The fundamental group $\pi_1(M)$ of a complete finite-volume hyperbolic manifold $M$ is residually finite.
\end{prop}
\begin{proof}
The group $\pi_1(M)$ is finitely generated (actually, finitely presented) because $M$ is homeomorphic to the interior of a compact manifold with boundary, so Proposition \ref{FG:RF:prop} applies.
\end{proof}

This algebraic fact has some nice geometric consequences. Recall that $\ell_1(M)$ is the length of the shortest closed geodesic in $M$.

\begin{cor}
For every $L>0$ there is a finite cover $\tilde M \to M$ with $\ell_1(\tilde M)>L$.
\end{cor}
\begin{proof}
There are finitely many hyperbolic elements $a_1,\ldots, a_k \in \pi_1(M)$ such that every closed geodesic in $M$ of length $<L$ is freely homotopic to some of these.
By residually finiteness there is a finite-index normal $H \triangleleft \pi_1(M)$ that does not contain any of the elements $a_1,\ldots, a_k$, and hence none of their conjugates. 

The finite-index $H$ determines a finite cover $\pi\colon \tilde M \to M$ with $\pi_*(\pi_1(\tilde M)) = H$. A closed geodesic in $\tilde M$ cannot be shorter than $L$ because its image in $M$ would be freely homotopic to one of $a_1,\ldots, a_k$.
\end{proof}

\begin{cor}
If $M$ is a closed hyperbolic manifold, for every $R>0$ there is a finite cover $\tilde M \to M$ with $\inj\tilde M > R$.
\end{cor}

We can summarise this by saying that every closed hyperbolic manifold has arbitrarily fat finite covers.

\subsection{Subgroup separability}
The geometric consequences of residually finiteness are quite remarkable: we now introduce a stronger algebraic notion that will lead to more geometric applications.

Let $G$ be a group. A subgroup $H<G$ is \emph{separable} if the intersection of all the finite-index subgroups of $G$ containing $H$ is $H$ itself. In other words, for every non-trivial element $a\in G\setminus H$, there is a finite-index subgroup $G'< G$ that contains $H$ but not $a$ and hence ``separates'' $H$ from $a$.\index{separable subgroup} 

By definition, the trivial subgroup $\{e\}$ in $G$ is separable if and only if $G$ is residually finite. Note that, for general $H<G$, we cannot require the subgroup $G'$ to be normal in $G$, because $H$ itself may not be normal in $G$, as opposite to $\{e\}$.
The following proposition furnishes some interesting examples.

\begin{prop} \label{Long:prop}
Every maximal abelian subgroup $H$ in a residually finite group $G$ is separable.
\end{prop}
\begin{proof}
More generally, we suppose that $H$ is maximal with respect to some word relation $f(h_1,\ldots, h_n)=e$. (In our case, $f(h_1, h_2) = [h_1, h_2]$.) Since $G$ is residually finite, there is a sequence $N_i \triangleleft G $ of finite-index normal subgroups with $\cap N_i = \{e\}$. 

We show that $H = \cap_i HN_i$, and we conclude since $HN_i > N_i$ has finite index in $G$. We have $H \subset \cap_i HN_i$. We now show that the elements of $\cap_i HN_i$ satisfy $f=e$, and by maximality of $H$ we get $H=\cap_i HN_i$.

For every $N_i$ and $h_1,\ldots, h_n \in H$, we have $f(h_1N_i, \ldots, h_nN_i) \subset  N_i$ since its projection in $G/_{N_i}$ is trivial. Therefore $\cap_i f(HN_i, \ldots, HN_i) = \{e\}$ and the elements of $\cap_i HN_i$ satisfy $f=e$.
\end{proof}

\begin{cor} \label{separable:cor}
Let $M$ be a complete finite-volume hyperbolic $n$-manifold. The following subgroups are separable:
\begin{itemize}
\item the subgroup $\langle \varphi \rangle = \matZ$ generated by a primitive hyperbolic transformation $\varphi$;
\item the subgroup $\pi_1(T)$ generated by a $(n-1)$-torus cusp section $T$.
\end{itemize}
\end{cor}
\begin{proof}
Both subgroups are maximal abelian, see Section \ref{that:commute:subsection}.
\end{proof}

We deduce a couple of geometric consequences from this algebraic fact.

\begin{cor}
Let $M$ be a complete finite-volume hyperbolic $n$-manifold. Every closed geodesic $\gamma$ in $M$ lifts to a closed geodesic in some finite covers of $M$ of arbitrarily large degree.
\end{cor}
\begin{proof}
We have $M= \matH^n/_\Gamma$. The closed geodesic $\gamma$ is obtained by projecting the axis $l$ of a primitive hyperbolic element $\varphi \in \Gamma$.
Since $\langle \varphi \rangle$ is separable, there are subgroups $H<\pi_1(M)$ of arbitrarily large index that contain $\langle \varphi \rangle$.
\end{proof}

We have discovered in particular that there are many closed hyperbolic manifolds of bounded injectivity radius with arbitrarily large volume.
We can also use separability to promote some primitive closed geodesics to simple ones on some finite coverings.

\begin{cor}
Let $M$ be a complete finite-volume hyperbolic $n$-manifold. Every primitive closed geodesic $\gamma$ in $M$ lifts to a simple closed geodesic in some finite cover of $M$.
\end{cor}
\begin{proof}
We have $M= \matH^n/_\Gamma$. The closed geodesic $\gamma$ is obtained by projecting the axis $l$ of a primitive hyperbolic element $\varphi \in \Gamma$. If $\gamma$ is not simple, it self-intersects in $k>0$ points, which lift to $k$ transverse intersections between $l$ and some translates $\varphi_1(l), \ldots, \varphi_k(l)$, for some $\varphi_i \in \Gamma$. Therefore the elements $\psi \in \Gamma$ such that $\psi(l)$ intersects $l$ transversely are those of the form $\psi = \varphi^h\varphi_i\varphi^l$ for some $i=1,\ldots, k$ and $h,l \in \matZ$.

Since $\langle \varphi \rangle$ is separable, there is a finite-index $H<\Gamma$ with $\varphi_1, \ldots, \varphi_k \not\in H$ and $\langle \varphi \rangle \subset H$. No $\varphi^h\varphi_i\varphi^l$ lies in $H$, hence for every $\psi\in H$ the lines $l$ and $\psi(l)$ coincide or are disjoint. The line $l$ projects to a simple geodesic in $\matH^n/_H$.
\end{proof}

Concerning cusps, we deduce analogously the following.

\begin{cor}
Let $M$ be a complete finite-volume hyperbolic $n$-manifold and $T\subset M$ be a $(n-1)$-torus cusp section. The cusp section $T$ lifts to a cusp section in some finite covers of $M$ with arbitrarily large degree.
\end{cor}
\begin{proof}
The subgroup $\pi_1(T)$ is separable, so there are subgroups $H<\pi_1(M)$ of arbitrarily large finite degree containing it.
\end{proof}

\section{The Bieberbach Theorem} \label{Bieberbach:section}
The Margulis Lemma for Lie groups is fairly general and has important applications also in the elliptic and flat geometries. The most important one is the Bieberbach Theorem.

\subsection{Elliptic manifolds}
Recall that every complete elliptic manifold is isometric to $S^n/_\Gamma$ for some finite group $\Gamma < \On(n) = \Iso(S^{n-1})$ acting freely. We now improve the Margulis Lemma \ref{G:lemma} for $G=\On(n)$ by promoting ``nilpotent'' to ``abelian''. To this purpose we need the following.

\begin{lemma} \label{backwards:lemma}
There is a neighbourhood $U$ of $I\in \On(n)$ such that for every $A\in \On(n)$ and every $B\in U$ we have
$$[A,[A,B]] = I\ \Longrightarrow\ [A,B] = I.$$
\end{lemma}
\begin{proof}
We prove the lemma for the bigger Lie group $U(n)$ consisting of all unitary matrices. An immediate computation shows that
$$[A,[A,B]]=I \ \Longleftrightarrow\ [A, BA^{-1}B^{-1}] = I.$$
Let $\matC^n = \oplus V_i$ be the decomposition into orthogonal eigenspaces for $A$ and $A^{-1}$ (with respect to the standard hermitian product of $\matC^n$). The decomposition for the conjugate $BA^{-1}B^{-1}$ is just $\matC^n = \oplus B(V_i)$. 

Let $U$ be the neighbourhood of $I$ consisting of all matrices that move every vector of an angle $< \frac \pi 2$. Pick $B\in U$. Since the $V_i$ are orthogonal we get $B(V_i) \cap V_j = \{0\}$ for all $i\neq j$. On the other hand, the endomorphisms $A$ and $BA^{-1}B^{-1}$ commute and hence have a basis of common eigenvectors: hence the only possibility is that $B(V_i) = V_i$ for all $i$. 

The restriction $A|_{V_i}$ is just $\lambda_iI$ and hence commutes with $B|_{V_i}$, for all $i$. Therefore $A$ and $B$ commute everywhere.
\end{proof}

\begin{cor} \label{U:O(n):cor}
There is a neighbourhood $U$ of $I\in \On(n)$ such that every finite subgroup $\Gamma < \On(n)$ generated by some elements in $U$ is abelian.
\end{cor}
\begin{proof}
We know from Margulis Lemma \ref{G:lemma} that there is a $U$ where every such $\Gamma$ is nilpotent. The previous lemma promotes $\Gamma$ to an abelian group, because $[A,[A, \ldots [A,B]\cdots]] = I$ implies after finitely many steps that $[A,B] = I$ for every generators $A,B\in \Gamma \cap U$.
\end{proof}

\begin{prop} \label{U:finite:index:prop}
Let $G$ be a compact Lie group and $U$ a neighbourhood of $e\in G$. There is a $N>0$ such that for every group $\Gamma < G$, the subgroup $\Gamma_U< \Gamma$ generated by $\Gamma \cap U$ has index at most $N$ in $\Gamma$.
\end{prop}
\begin{proof}
Let $W\subset U$ be a smaller neighbourhood such that $W^{-1}=W$ and $W^2\subset U$ and set $N = \Vol (G) / \Vol (W)$ using the Haar measure for $G$.

We conclude by showing that if $g\Gamma_U$ and $g'\Gamma_U$ are distinct cosets of $\Gamma_U$ in $\Gamma$, then $gW \cap g'W = \emptyset$. This implies that $\Gamma_U$ has index at most $N$ in $\Gamma$.

Indeed, if $gW \cap g'W \neq \emptyset$ there are $w,w'\in W$ such that $gw = g'w'$ which implies that $g^{-1}g' = w(w')^{-1}\in W^2\subset U$ and hence $g^{-1}g' \in \Gamma_U$.
\end{proof}

We obtain a fairly interesting corollary about finite subgroups of $\On(n)$.

\begin{cor} \label{n:N:cor}
For every $n$ there is a $N>0$ such that every finite subgroup of $\On(n)$ contains an abelian subgroup of index at most $N$.
\end{cor}

\subsection{Isometries of Euclidean space}
We now turn to flat manifolds. Recall that a complete flat manifold is isometric to $\matR^n/_\Gamma$ for some discrete group $\Gamma < \Iso(\matR^n)$ acting freely. Every isometry $g$ of $\matR^n$ can be written uniquely as
$$g\colon x \longmapsto Ax +b$$
for some $A\in \On(n)$ and $b\in \matR^n$. The \emph{rotational} and \emph{translational part} of $g$ are $A$ and $b$ respectively. Let $\Fix(A)$ be the fixed points of $x\mapsto Ax$.

\begin{prop} \label{basic:isometries:Rn:prop}
Let $Ax+b$ be an isometry of $\matR^n$. Then:
\begin{itemize}
\item its inverse is $A^{-1}x-A^{-1}b$,
\item if it acts freely then $\Fix(A)\neq \{0\}$,
\item $\exists$ a translation of $\matR^n$ that conjugates it to $Ax+b'$ with $b'\in \Fix(A)$.
\end{itemize}
\end{prop}
\begin{proof}
If it acts freely then $Ax+b=x$ has no solution $x\in\matR^n$. By rewriting the equation as $(A-I)x = -b$ we see that $A-I$ is not surjective, hence not injective, hence $\Fix(A)$ is non-trivial.

A translation $x+d$ conjugates $Ax+b$ into 
$$\big(A(x+d)+b\big)-d = Ax + (A-I)d + b.$$
Since $A$ is orthogonal, we get $\Img (A-I) = \ker (A-I)^\perp = \Fix(A)^\perp$ and hence there is a $d$ such that $b'=(A-I)d+b\in \Fix(A)$.
\end{proof}

\begin{ex} \label{Euclidean:commutators:ex}
The commutator of two isometries is
$$[Ax+b, Cx+d] = [A,C]x + A(I-C)A^{-1}b + AC(I-A^{-1})C^{-1}d.$$
In particular we get
$$[Ax+b,x+d] = x + (A-I)d.$$
\end{ex}

\subsection{Discrete groups}
The homomorphism $r\colon \Iso(\matR^n)\to \On(n)$ that sends every isometry to its rotational part induces an exact sequence
$$0 \longrightarrow \matR^n \longrightarrow \Iso(\matR^n) \stackrel{r}\longrightarrow \On(n) \longrightarrow 0$$
where we indicate by $\matR^n$ the group of translations of $\matR^n$. 

Let $\Gamma < \Iso(\matR^n)$ be a discrete group. We get an exact sequence
$$0 \longrightarrow H \longrightarrow \Gamma \longrightarrow r(\Gamma) \longrightarrow 0$$
where $H \triangleleft \Gamma$ is the \emph{translation subgroup} of $\Gamma$. The subgroup $r(\Gamma)<\On(n)$ is not necessarily discrete, as the following example shows.

\begin{example} \label{rototranslation:example}
A \emph{rototranslation} in $\matR^3$ is a rotation of some angle $\theta$ along an axis $r$ composed with a translation of some distance $t>0$ in the direction of $r$. A rototranslation generates a discrete group $\Gamma = \matZ$ acting freely on $\matR^3$, whose quotient $\matR^3/_\Gamma$ is diffeomorphic to $\matR^2\times S^1$. If $\theta$ is not commensurable with $\pi$ the image $r(\Gamma)$ is not discrete and forms a dense subset of the circle in $\On(3)$ of all rotations along $r$.
\end{example}

We now extend Corollary \ref{n:N:cor} to this context.

\begin{teo} \label{discrete:virtual:abelian:teo}
For every $n$ there is a $N>0$ such that every discrete subgroup of $\Iso(\matR^n)$ contains an abelian subgroup of index at most $N$.
\end{teo}
\begin{proof}
We pick a small neighbourhood $U \subset \On(n)$ of $I$, such that
the commutator is contracting in $U$ (see the proof of the Margulis Lemma \ref{G:lemma}), it is symmetric ($U=U^{-1}$) and $A-I$ contracts vectors uniformly for all $A\in U$, more precisely $\|(A-I)v\| < \frac 14\|v\|$ for all $v\neq 0$.

Let $\Gamma < \Iso(\matR^n)$ be discrete and $r(\Gamma)_U$ be the group generated by $r(\Gamma)\cap U$. By Proposition \ref{U:finite:index:prop} the index of $r(\Gamma)_U$ in $r(\Gamma)$ is bounded by some $N$ depending only on $U$. Therefore its counterimage in $\Gamma$
$$\Gamma^* = r^{-1}(r(\Gamma)_U) \cap \Gamma$$ 
has index in $\Gamma$ bounded by $N$. The group $\Gamma^*$ is the subgroup of $\Gamma$ consisting of all $Ax+b$ with $A\in r(\Gamma)_U$. 
It remains to prove that $\Gamma^*$ is abelian.

Let $A_1x+b_1$ and $A_2x +b_2$ be two elements in $\Gamma^*$. Define
\begin{equation*} 
A_{i+1}x+b_{i+1} = [A_1x+b_1, A_ix+b_i]
\end{equation*}
for all $i\geqslant 2$. Exercise \ref{Euclidean:commutators:ex} gives
\begin{equation} \label{commutator:eqn}
A_{i+1}x+b_{i+1} = [A_1,A_i]x + A_1(I-A_i)A_1^{-1}b_1 + A_1A_i(I-A_1^{-1})A_i^{-1}b_i
\end{equation}
and hence
\begin{align*}
A_{i+1} & = [A_1,A_i], \\
 b_{i+1} & = A_1(I-A_i)A_1^{-1}b_1 + A_1A_i(I-A_1^{-1})A_i^{-1}b_i.
\end{align*}
Suppose $A_1, A_2 \in U$. Since the commutator is contracting in $U$ and $A_i-I$ contract vectors uniformly, we get $A_{i} \to I$ and $b_i\to 0$ as $i\to \infty$. Therefore $A_ix+b_i$ tends to the identity, and since $\Gamma^*$ is discrete it \emph{is} the identity for all $i$ bigger or equal than some $i_0$. In particular $A_{i_0} = I$ and Lemma \ref{backwards:lemma} used backwards gives $A_3=[A_1,A_2]=I$: hence $A_1$ and $A_2$ commute. 

Since $r(\Gamma^*)$ is generated by $r(\Gamma^*)\cap U$,
we deduce that $r(\Gamma^*)$ is abelian. Therefore (\ref{commutator:eqn}) may be restated as
\begin{equation} \label{simplified:commutator:eqn}
A_{i+1}x+b_{i+1} = x + (I-A_i)b_1 + (A_1 - I)b_i.
\end{equation}
We now consider the case $A_1 \in U$ and $A_2=I$. We get
$$A_{i+1}x+b_{i+1} = x + (A_1-I)^{i-1} b_2.$$
Since $A_1\in U$ we get $(A_1-I)^i b_2 \to 0$ and hence $(A_1-I)^ib_2=0$ for some $i$, which gives $(A_1-I)b_2=0$ since $A_1$ is diagonalisable. Thus $b_2 \in \Fix (A_1)$.

We have proved that if $\Gamma^*$ contains a translation $x+b$, then $b\in \Fix(A)$ for all $A\in r(\Gamma^*)\cap U$. Since these $A$ generate $r(\Gamma^*)$, the vector $b$ belongs to 
$$W = \Fix(r(\Gamma^*)) = \big\{x\ \big|\ Ax = x\ \forall A \in r(\Gamma^*)\big\}.$$
Pick now two arbitrary elements $Ax+b$ and $Cx+d$ in $\Gamma^*$.
By (\ref{simplified:commutator:eqn}) we get
$$[Ax+b, Cx+d] = x+(I-C)b+(A-I)d.$$
By what just said $(I-C)b+(A-I)d \in W$. On the other hand 
$$\Img (I-C) = \ker(I-C)^\perp = \Fix(C)^\perp \subset W^\perp$$ 
and hence $(I-C)b\in W^\perp$, and analogously $(A-I)d \in W^\perp$. We deduce that $(I-C)b+(A-I)d \in W \cap W^\perp$ is trivial and all elements in $\Gamma^*$ commute.
\end{proof}


\subsection{Crystallographic groups} \label{crystallographic:subsection}
A \emph{crystallographic group} is a discrete subgroup $\Gamma <\Iso(\matR^n)$ with compact quotient $\matR^n/_\Gamma$.\index{group!crystallographic group}

\begin{prop} \label{finite:image:r:prop}
The image $r(\Gamma)$ of a crystallographic group $\Gamma<\Iso(\matR^n)$ is finite. 
\end{prop}
\begin{proof}
By Theorem \ref{discrete:virtual:abelian:teo} we may suppose that $\Gamma$ is abelian. We now prove that $\Gamma$ abelian implies that $r(\Gamma)$ is trivial, \emph{i.e.}~all elements in $\Gamma$ are translations.

Suppose that $\Gamma$ contains a non-translation $Ax +b$. We conjugate $\Gamma$ by a translation as in Proposition \ref{basic:isometries:Rn:prop} to get $b\in \Fix(A)$. Pick another isometry $Cx+d$ in $\Gamma$. The commutator
$$[Ax +b, Cx +d] = x+(A-I)d-(C-I)b$$
is trivial. Since $A$ and $C$ commute, we get $(C-I)b \in \Fix(A)$, hence $(A-I)
d\in \Fix (A)$ and finally $d\in \Fix(A)$.

We have proved that $d\in \Fix(A)\subsetneq \matR^n$ for all elements $Cx+d$ in $\Gamma$. Therefore the $\Gamma$-orbit of $0\in\matR^n$ is contained in $\Fix(A)$. But the compactness of $\matR^n/_\Gamma$ implies that there is a compact fundamental domain, and hence a $R>0$ such that every point in $\matR^n$ is $R$-close to any fixed orbit: a contradiction since $\Fix(A)$ is a proper vector subspace of $\matR^n$.
\end{proof}

\begin{cor} \label{crystal:cor}
Every crystallographic group has a finite-index translation subgroup isomorphic to $\matZ^n$.
\end{cor}

Recall from Section \ref{piatte:subsection} that a \emph{flat torus} is a $n$-torus $\matR^n/_\Gamma$ where $\Gamma$ is a lattice, \emph{i.e.}~a discrete group isomorphic to $\matZ^n$ that spans $\matR^n$ as a vector space.\index{flat torus}\index{Bieberbach Theorem}

\begin{cor}[Bieberbach's Theorem] \label{Bieberbach:cor}
Every closed flat $n$-manifold is finitely covered by a flat torus.
\end{cor}

We conclude the discussion by noting that there are no ``cusps'' in Euclidean geometry.

\begin{prop} \label{finite:volume:flat:prop}
Every finite-volume complete flat manifold is closed.
\end{prop}
\begin{proof}
Let $M=\matR^n/_\Gamma$ be a finite-volume flat manifold. Up to taking finite indexes, the group $\Gamma$ is abelian by Theorem \ref{discrete:virtual:abelian:teo}. If all the elements are translations, we are done: finite-volume easily implies that $\Gamma = \matZ^n$, thus $M$ is closed. 

Suppose some element $Ax + b$ is not a translation: every element of $\Gamma$ commutes with it and hence preserves the space $\Fix (A)$ that has dimension $k<n$; hence $M$ is isometric to $\Fix (A)/_\Gamma \times \matR^{n-k}$ and has infinite volume.
\end{proof}

\subsection{References}
Most of the material introduced in this chapter is standard and can be found in Benedetti -- Petronio \cite{BP}, Ratcliffe \cite{R}, and of course in Thurston's notes \cite{Th}. We have also consulted Thurston's book \cite{Th_book} for the part on chrystallographic groups. Proposition \ref{Long:prop} was taken from a paper of Long \cite{L}.


